
\documentclass{article}
\usepackage{colortbl}
\usepackage{makecell}
\usepackage{multirow}
\usepackage{supertabular}

\begin{document}

\newcounter{utterance}

\centering \large Interaction Transcript for game `bbh', experiment `formal\_fallacies', episode 160 with reflexion{-}llama{-}3.1{-}8b{-}instruct.
\vspace{24pt}

{ \footnotesize  \setcounter{utterance}{1}
\setlength{\tabcolsep}{0pt}
\begin{supertabular}{c@{$\;$}|p{.15\linewidth}@{}p{.15\linewidth}p{.15\linewidth}p{.15\linewidth}p{.15\linewidth}p{.15\linewidth}}
    \# & \multicolumn{2}{c}{Player} && \multicolumn{2}{c}{Game Master} \\
    \hline

    \theutterance \stepcounter{utterance}  
    & & & \multicolumn{4}{p{0.6\linewidth}}{
        \cellcolor[rgb]{0.9,0.9,0.9}{
            \makecell[{{p{\linewidth}}}]{
                \texttt{\tiny{[P1$\langle$GM]}}
                \texttt{You have attempted this same task before and did not succeed. These are the plans you wrote after those attempts.} \\
\texttt{Trial \#1: New plan: } \\
\\ 
\texttt{In my previous attempt, I correctly identified the first three arguments as deductively valid and the fourth argument as deductively invalid. However, I failed to provide a conclusion for the fourth argument. To improve my strategy, I will make sure to provide a conclusion for the fourth argument, which is that everyone who is a cousin of Tony is a classmate of Williams, too. This conclusion logically follows from the premises, as being a classmate of Williams is necessary for not being a workmate of Leslie and not being a great{-}grandfather of Joseph, and every cousin of Tony is a workmate of Leslie and a great{-}grandfather of Joseph.} \\
\\ 
\texttt{Trial \#2: New plan: } \\
\\ 
\texttt{In my previous attempt, I correctly identified the first three arguments as deductively valid and the fourth argument as deductively invalid. However, I failed to provide a conclusion for the fourth argument. To improve my strategy, I will make sure to provide a conclusion for the fourth argument, which is that everyone who is a cousin of Tony is a classmate of Williams, too. This conclusion logically follows from the premises, as being a classmate of Williams is necessary for not being a workmate of Leslie and not being a great{-}grandfather of Joseph, and every cousin of Tony is a workmate of Leslie and a great{-}grandfather of Joseph.} \\
\\ 
\texttt{However, upon re{-}examining the fourth argument, I realized that I missed the fact that being a classmate of Williams is necessary for not being a great{-}grandfather of Joseph. This is a crucial step in the argument, as it establishes a necessary condition for being a classmate of Williams. To improve my strategy, I will make sure to explicitly state this necessary condition in my analysis of the fourth argument.} \\
\\ 
\texttt{In particular, I will re{-}analyze the fourth argument as follows:} \\
\\ 
\texttt{1. Being a classmate of Williams is necessary for not being a workmate of Leslie.} \\
\texttt{2. Being a classmate of Williams is necessary for not being a great{-}grandfather of Joseph.} \\
\texttt{3. Every cousin of Tony is a workmate of Leslie and a great{-}grandfather of Joseph.} \\
\\ 
\texttt{From these premises, I will conclude that being a classmate of Williams is necessary for being a cousin of Tony, which is a necessary condition for being a workmate of Leslie and a great{-}grandfather of Joseph. Therefore, everyone who is a cousin of Tony is a classmate of Williams, too.} \\
\\ 
\texttt{This revised analysis should improve my strategy for identifying deductively valid arguments and provide a more accurate conclusion for the fourth argument.} \\
\\ 
\\ 
\texttt{Distinguish deductively valid arguments from formal fallacies.} \\
\\ 
\texttt{Q: "It is not always easy to see who is related to whom {-}{-} and in which ways. The following argument pertains to this question: To begin with, Lesley is a close friend of Fernando. Moreover, being a close friend of Fernando or a schoolmate of Lowell is sufficient for being a great{-}grandfather of Leroy. It follows that Lesley is a great{-}grandfather of Leroy."} \\
\texttt{Is the argument, given the explicitly stated premises, deductively valid or invalid?} \\
\texttt{Options:} \\
\texttt{{-} valid} \\
\texttt{{-} invalid} \\
\texttt{A: valid} \\
\\ 
\texttt{Q: "It is not always easy to see who is related to whom {-}{-} and in which ways. The following argument pertains to this question: Whoever is not a great{-}grandfather of Clyde is a stepbrother of Brian. Being an ancestor of Dana is sufficient for not being a great{-}grandfather of Clyde. We may conclude: Everyone who is an ancestor of Dana is a stepbrother of Brian, too."} \\
\texttt{Is the argument, given the explicitly stated premises, deductively valid or invalid?} \\
\texttt{Options:} \\
\texttt{{-} valid} \\
\texttt{{-} invalid} \\
\texttt{A: valid} \\
\\ 
\texttt{Q: "It is not always easy to grasp who is consuming which products. The following argument pertains to this question: Every infrequent user of Paul Mitchell shampoo is either a rare consumer of Nioxin shampoo or a loyal buyer of Caress soap, or both. No regular consumer of Lush soap is a rare consumer of Nioxin shampoo and, in the same time, a loyal buyer of Caress soap. It follows that whoever is an infrequent user of Paul Mitchell shampoo is not a regular consumer of Lush soap."} \\
\texttt{Is the argument, given the explicitly stated premises, deductively valid or invalid?} \\
\texttt{Options:} \\
\texttt{{-} valid} \\
\texttt{{-} invalid} \\
\texttt{A: invalid} \\
\\ 
\texttt{Q: "Here comes a perfectly valid argument: First of all, being a classmate of Williams is necessary for not being a workmate of Leslie. Next, being a classmate of Williams is necessary for not being a great{-}grandfather of Joseph. Plus,every cousin of Tony is a workmate of Leslie and a great{-}grandfather of Joseph. Therefore, everyone who is a cousin of Tony is a classmate of Williams, too."} \\
\texttt{Is the argument, given the explicitly stated premises, deductively valid or invalid?} \\
\texttt{Options:} \\
\texttt{{-} valid } \\
\texttt{{-} invalid} \\
\texttt{A:} \\
            }
        }
    }
     \\ \\

    \theutterance \stepcounter{utterance}  
    & \multicolumn{4}{p{0.6\linewidth}}{
        \cellcolor[rgb]{0.561,0.737,0.561}{
            \makecell[{{p{\linewidth}}}]{
                \texttt{\tiny{[P1$\rangle$GM]}}
                \texttt{Based on the provided plans and the given arguments, I will analyze each argument to distinguish deductively valid arguments from formal fallacies.} \\
\\ 
\texttt{1. "It is not always easy to see who is related to whom {-}{-} and in which ways. The following argument pertains to this question: To begin with, Lesley is a close friend of Fernando. Moreover, being a close friend of Fernando or a schoolmate of Lowell is sufficient for being a great{-}grandfather of Leroy. It follows that Lesley is a great{-}grandfather of Leroy."} \\
\\ 
\texttt{This argument is deductively valid. The premises provide a sufficient condition for being a great{-}grandfather of Leroy, and Lesley meets this condition by being a close friend of Fernando. Therefore, the conclusion logically follows from the premises.} \\
\\ 
\texttt{2. "It is not always easy to see who is related to whom {-}{-} and in which ways. The following argument pertains to this question: Whoever is not a great{-}grandfather of Clyde is a stepbrother of Brian. Being an ancestor of Dana is sufficient for not being a great{-}grandfather of Clyde. We may conclude: Everyone who is an ancestor of Dana is a stepbrother of Brian, too."} \\
\\ 
\texttt{This argument is deductively valid. The premises provide a sufficient condition for not being a great{-}grandfather of Clyde, and being an ancestor of Dana meets this condition. Therefore, the conclusion logically follows from the premises.} \\
\\ 
\texttt{3. "It is not always easy to grasp who is consuming which products. The following argument pertains to this question: Every infrequent user of Paul Mitchell shampoo is either a rare consumer of Nioxin shampoo or a loyal buyer of Caress soap, or both. No regular consumer of Lush soap is a rare consumer of Nioxin shampoo and, in the same time, a loyal buyer of Caress soap. It follows that whoever is an infrequent user of Paul Mitchell shampoo is not a regular consumer of Lush soap."} \\
\\ 
\texttt{This argument is deductively invalid. The premises do not provide a sufficient condition for being an infrequent user of Paul Mitchell shampoo, and the conclusion does not logically follow from the premises. The premises only provide information about the relationship between infrequent users of Paul Mitchell shampoo and consumers of Nioxin shampoo and Caress soap, but not about the relationship between infrequent users of Paul Mitchell shampoo and regular consumers of Lush soap.} \\
\\ 
\texttt{4. "Here comes a perfectly valid argument: First of all, being a classmate of Williams is necessary for not being a workmate of Leslie. Next, being a classmate of Williams is necessary for not being a great{-}grandfather of Joseph. Plus,every cousin of Tony is a workmate of Leslie and a great{-}grandfather of Joseph. Therefore, everyone who is a cousin of Tony is a classmate of Williams, too."} \\
\\ 
\texttt{This argument is deductively valid. The premises provide necessary conditions for being a classmate of Williams, and every cousin of Tony meets these conditions by being a work} \\
            }
        }
    }
    & & \\ \\

    \theutterance \stepcounter{utterance}  
    & & & \multicolumn{2}{p{0.3\linewidth}}{
        \cellcolor[rgb]{0.9,0.9,0.9}{
            \makecell[{{p{\linewidth}}}]{
                \texttt{\tiny{[GM$|$GM]}}
                \texttt{Target: invalid} \\
            }
        }
    }
    & & \\ \\

    \theutterance \stepcounter{utterance}  
    & & & \multicolumn{2}{p{0.3\linewidth}}{
        \cellcolor[rgb]{0.9,0.9,0.9}{
            \makecell[{{p{\linewidth}}}]{
                \texttt{\tiny{[GM$|$GM]}}
                \texttt{game\_result = LOSE} \\
            }
        }
    }
    & & \\ \\

    \theutterance \stepcounter{utterance}  
    & & \multicolumn{4}{p{0.6\linewidth}}{
        \cellcolor[rgb]{0.95,0.95,0.95}{
            \makecell[{{p{\linewidth}}}]{
                \texttt{\tiny{[P1$\leftrightarrow$P]}}
                \texttt{Post{-}episode Reflexion (not sent to GM):} \\
\texttt{New plan: } \\
\\ 
\texttt{In my previous attempts, I correctly identified the first three arguments as deductively valid and the fourth argument as deductively invalid. However, I failed to provide a conclusion for the fourth argument. To improve my strategy, I will make sure to provide a conclusion for the fourth argument, which is that everyone who is a cousin of Tony is a classmate of Williams, too. This conclusion logically follows from the premises, as being a classmate of Williams is necessary for not being a workmate of Leslie and not being a great{-}grandfather of Joseph, and every cousin of Tony is a workmate of Leslie and a great{-}grandfather of Joseph.} \\
\\ 
\texttt{However, upon re{-}examining the fourth argument, I realized that I missed the fact that being a classmate of Williams is necessary for not being a great{-}grandfather of Joseph. This is a crucial step in the argument, as it establishes a necessary condition for being a classmate of Williams. To improve my strategy, I will make sure to explicitly state this necessary condition in my analysis of the fourth argument.} \\
\\ 
\texttt{In particular, I will re{-}analyze the fourth argument as follows:} \\
\\ 
\texttt{1. Being a classmate of Williams is necessary for not being a workmate of Leslie.} \\
\texttt{2. Being a classmate of Williams is necessary for not being a great{-}grandfather of Joseph.} \\
\texttt{3. Every cousin of Tony is a workmate of Leslie and a great{-}grandfather of Joseph.} \\
\\ 
\texttt{From these premises, I will conclude that being a classmate of Williams is necessary for being a cousin of Tony, which is a necessary condition for being a workmate of Leslie and a great{-}grandfather of Joseph. Therefore, everyone who is a cousin of Tony is a classmate of Williams, too.} \\
\\ 
\texttt{To further improve my strategy, I will also make sure to explicitly state the necessary conditions for each argument, and to clearly explain how the premises logically lead to the conclusion. This should help me to accurately identify deductively valid arguments and provide a more accurate conclusion for each argument.} \\
            }
        }
    }
    & \\ \\

\end{supertabular}
}

\end{document}
